% META: {"id":"TCOMPL-LOG-CO-001","chapitre":"TCOMPL-LOGARITHME-HISTORIQUE","type_objet":"corrige","exercice_ref":"TCOMPL-LOG-EX-001","capacites_codes":["C4"],"sources_inspiration":[],"status":"approved","origin":"generated","status_history":["generated","audited","accepted","approved"],"release_acceptance":"RELEASE_OWNER_BATCH_ACCEPTANCE","acceptance_closure_digest":"sha256:9b3ccf9a81c5520fbb7e03b7b2d3e7bf2057a02834b6d908bf3be5f49f3b553f","acceptance_receipt_digest":"sha256:4fad86f0282e0fa4e0419cca1ad6379ae7af5a280c04a4a90880f90a1c044242"}
% BEGIN-VERIFY
% from sympy import ln, simplify, Rational
% A = ln(12) - ln(3) - ln(2)
% assert simplify(A - ln(2)) == 0
% B = ln(Rational(1,5)) + ln(5)
% assert simplify(B) == 0
% END-VERIFY
\begin{corrige}{TCOMPL-LOG-EX-001}
\textbf{1.} On regroupe les deux termes soustraits : $\ln(3)+\ln(2) = \ln(3\times2) = \ln(6)$. Donc :
\[
  A = \ln(12) - \ln(6) = \ln\!\left(\frac{12}{6}\right) = \ln(2).
\]

\textbf{2.} $\ln\!\left(\dfrac{1}{5}\right) = -\ln(5)$ (propriété de l'inverse), donc :
\[
  B = -\ln(5) + \ln(5) = 0.
\]
\end{corrige}
